\chapter{Existence of a prime factorisation}
\label{ch:factorization}

\section{The shape of the statement}

The source states factorisation in the classical exponent form,
$n = \prod_{i=1}^{k} p_i^{e_i}$ with the $p_i$ \emph{distinct} and $e_i > 0$.
This chapter proves it in \emph{both} shapes, and the distinction is one of
working convenience only.

The induction is run on a \textbf{list of primes whose product is $n$}, with
multiplicity recorded by repetition rather than by an exponent
(Theorem~\ref{thm:exists_factorization}). That is the form the argument wants:
its inductive step is literally list concatenation, whereas merging two
exponent-indexed families would require adding exponents at shared primes before
the induction has any grouping machinery available.

The source's exponent form is then recovered in \S\ref{sec:exponent_form} by
grouping the list by value: the distinct primes are the \emph{distinct} values
occurring in the list, and each exponent is the multiplicity with which its prime
occurs. Both statements are therefore project theorems; neither is asserted as
``equivalent up to an unproved translation''.

Only Chapter~\ref{ch:infinitude} consumes this chapter, and it needs neither
distinctness nor exponents --- it needs a single prime divisor
(Corollary~\ref{cor:exists_prime_dvd}). The exponent form is proved because the
source names it as a theorem, not because anything downstream calls it.

Uniqueness of the factorisation is deliberately \emph{not} part of this project.
The source itself flags it as unnecessary, and nothing in
Chapter~\ref{ch:infinitude} appeals to it.

\section{Reducibility produces two non-unit factors}

\begin{lemma}[unfolding the failure of irreducibility]
  \leanok
  \label{lem:exists_factor_of_not_irreducible}
  \lean{EuclidPrimes.exists_factor_of_not_isIrreducible}
  \uses{def:is_irreducible}
  Let $n \in \NN$ with $1 < n$, and suppose $\Irr{n}$ fails. Then there exist
  $a, b \in \NN$ with $n = a b$, $a \neq 1$ and $b \neq 1$.
\end{lemma}
\begin{proof}

\leanok
  $\Irr{n}$ is the conjunction of $1 < n$ with the factorisation clause. Since
  $1 < n$ holds by hypothesis, it is the factorisation clause that fails: there
  are $a, b$ with $n = a b$, $n \neq a$ and $n \neq b$.

  It remains to convert ``$n \neq a$ and $n \neq b$'' into ``$a \neq 1$ and
  $b \neq 1$''. If $a = 1$ then $n = 1 \cdot b = b$, contradicting $n \neq b$;
  so $a \neq 1$. If $b = 1$ then $n = a \cdot 1 = a$, contradicting $n \neq a$;
  so $b \neq 1$.
\end{proof}

\section{The factorisation theorem}

\begin{theorem}[existence of a prime factorisation]
  \leanok
  \label{thm:exists_factorization}
  \lean{EuclidPrimes.exists_factorization}
  \uses{def:is_prime, cor:prime_iff_irreducible, lem:exists_factor_of_not_irreducible, lem:one_lt_factor, lem:factor_lt}
  % SOURCE: [euclid-notes], the first `theroem:` (read from references/euclid-infinitude-of-primes.md)
  % SOURCE QUOTE: "theroem: For any natural number n > 1, there exists (p_1, e_1), ... (p_k, e_k) with p_i distict primes and e_i > 0 forall i
  % such that n = \Product_{1}^{k} p_i ^ e_i"
  \textit{Source: euclid-infinitude-of-primes.md, first \texttt{theroem:} (restated in list form; see \S1).}
  For every $n \in \NN$ with $1 < n$ there is a finite list
  $\ell = [q_1, \dots, q_k]$ of natural numbers such that every $q_i$ satisfies
  $\Prm{q_i}$ and
  \[
    \prod_{i=1}^{k} q_i = n .
  \]
\end{theorem}
% SOURCE QUOTE PROOF: "proof: induction on n
% base case: n = 2, immediate
% induction case, suppose true for 2 ... n
% If n + 1 is prime, then were done. Otherwise n + 1 is irreducible, so (n + 1) = a * b with niether a nor b eqaual to 1.
% Then a < n + 1 and b < n + 1, so we can apply the induction hypothesis to both a and b. Combine their factorizations to get the result."
\begin{proof}
  \uses{cor:prime_iff_irreducible, lem:exists_factor_of_not_irreducible, lem:one_lt_factor, lem:factor_lt}
  \leanok
  Two remarks on the source proof before we begin. First, its case split reads
  ``\texttt{If n + 1 is prime, then were done. Otherwise n + 1 is irreducible}'',
  which inverts the intended dichotomy: the ``otherwise'' branch is the one where
  $n+1$ is \emph{not} irreducible. We split on irreducibility and use
  Corollary~\ref{cor:prime_iff_irreducible} to move to primality. Second, the
  source's ``base case $n = 2$'' plus ``suppose true for $2 \dots n$'' is a
  strong induction; we run it as such, with no separate base case --- the case
  analysis below already covers $n = 2$, since $2$ is irreducible.

  We argue by strong induction on $n$: assume the statement for every $m < n$
  with $1 < m$, and prove it for $n$ with $1 < n$. Split on whether $\Irr{n}$
  holds.

  \paragraph{Case $\Irr{n}$.} By Corollary~\ref{cor:prime_iff_irreducible},
  $\Prm{n}$. Take $\ell := [n]$. Every member of $\ell$ is prime, and its product
  is $n$.

  \paragraph{Case $\lnot\,\Irr{n}$.} By
  Lemma~\ref{lem:exists_factor_of_not_irreducible} there are $a, b$ with
  $n = a b$, $a \neq 1$, $b \neq 1$. Then:
  \begin{itemize}
    \item $1 < a$ and $1 < b$, by Lemma~\ref{lem:one_lt_factor} applied to
      $n = a b$ and to $n = b a$.
    \item $a < n$ by Lemma~\ref{lem:factor_lt} applied to $n = a b$, and $b < n$
      by the same lemma applied to $n = b a$.
  \end{itemize}
  So the induction hypothesis applies to both $a$ and $b$, yielding lists
  $\ell_a$ and $\ell_b$ of primes with $\prod \ell_a = a$ and
  $\prod \ell_b = b$. Take $\ell := \ell_a \mathbin{+\!\!+} \ell_b$. A member of
  $\ell$ lies in $\ell_a$ or in $\ell_b$, hence is prime; and the product of a
  concatenation is the product of the products, so
  $\prod \ell = a b = n$.
\end{proof}

\section{Extracting a single prime divisor}

Chapter~\ref{ch:infinitude} does not need the whole factorisation, only that
every $n > 1$ has \emph{some} prime divisor. That is the theorem plus one
observation: a list whose product exceeds $1$ cannot be empty.

\begin{corollary}[every $n > 1$ has a prime divisor]
  \leanok
  \label{cor:exists_prime_dvd}
  \lean{EuclidPrimes.exists_isPrime_dvd}
  \uses{def:is_prime, thm:exists_factorization}
  For every $n \in \NN$ with $1 < n$ there exists $p$ with $\Prm{p}$ and
  $p \mid n$.
\end{corollary}
\begin{proof}
  \uses{thm:exists_factorization}
  \leanok
  Take $\ell$ as in Theorem~\ref{thm:exists_factorization}, so all members of
  $\ell$ are prime and $\prod \ell = n$. The list $\ell$ is non-empty: the
  product of the empty list is $1$, whereas $\prod \ell = n > 1$. Let $p$ be its
  first member. Then $\Prm{p}$ because $p \in \ell$, and $p \mid \prod \ell = n$
  because any member of a list divides that list's product.
\end{proof}

\section{Recovering the source's exponent form}
\label{sec:exponent_form}

It remains to produce the statement the source actually names: a factorisation
indexed by \emph{distinct} primes carrying strictly positive exponents. Nothing
downstream needs it, but it is one of the source's theorems, so it is proved
here rather than left as an informal remark about
Theorem~\ref{thm:exists_factorization}.

Distinctness is expressed by indexing the product over a finite \emph{set} of
primes: a set has no repeated elements, so ``the $p_i$ are distinct'' becomes a
property of the index rather than an extra hypothesis to carry. The exponents are
given by a function $e$ on $\NN$, constrained only at the primes that actually
occur; its values elsewhere are irrelevant to the product.

\begin{theorem}[existence of a prime factorisation, exponent form]
  \leanok
  \label{thm:exists_factorization_pow}
  \lean{EuclidPrimes.exists_factorization_pow}
  \uses{def:is_prime, thm:exists_factorization}
  % SOURCE: [euclid-notes], the first `theroem:` (read from references/euclid-infinitude-of-primes.md)
  % SOURCE QUOTE: "theroem: For any natural number n > 1, there exists (p_1, e_1), ... (p_k, e_k) with p_i distict primes and e_i > 0 forall i
  % such that n = \Product_{1}^{k} p_i ^ e_i"
  \textit{Source: euclid-infinitude-of-primes.md, first \texttt{theroem:} (the
  source's own exponent form; the list form is Theorem~\ref{thm:exists_factorization}).}
  For every $n \in \NN$ with $1 < n$ there exist a finite set $S \subseteq \NN$
  and a function $e : \NN \to \NN$ such that
  \begin{itemize}
    \item every $p \in S$ satisfies $\Prm{p}$,
    \item $e(p) > 0$ for every $p \in S$, and
    \item $\displaystyle\prod_{p \in S} p^{\,e(p)} = n$.
  \end{itemize}
\end{theorem}
\begin{proof}
  \uses{thm:exists_factorization}
  \leanok
  Let $\ell$ be a list of primes with $\prod \ell = n$, as supplied by
  Theorem~\ref{thm:exists_factorization}, and regard $\ell$ as a finite
  multiset --- the induction produced it in some order, but the order carries no
  information here.

  Take $S$ to be the set of \emph{distinct} values occurring in $\ell$, and for
  each $p$ let $e(p)$ be the multiplicity of $p$ in $\ell$, i.e.\ the number of
  positions of $\ell$ at which $p$ occurs. The three claims are then:

  \begin{itemize}
    \item \emph{Every element of $S$ is prime.} An element of $S$ occurs in
      $\ell$, and every member of $\ell$ is prime.
    \item \emph{Exponents are positive on $S$.} A value occurs in $\ell$ exactly
      when its multiplicity is non-zero, so $p \in S$ gives $e(p) > 0$.
    \item \emph{The product is $n$.} Collecting equal factors of a finite
      product in a commutative monoid, the product of a multiset equals the
      product, over its distinct values, of each value raised to its
      multiplicity. Applied to $\ell$ this reads
      $\prod_{p \in S} p^{\,e(p)} = \prod \ell = n$.
  \end{itemize}

  Note that $S$ is automatically non-empty, since an empty index set would give
  product $1 \neq n$; the statement does not need to say so separately.
\end{proof}

\section{Uniqueness of the factorisation: deliberately out of scope}
\label{sec:uniqueness_out_of_scope}

The source continues past existence with two further named results --- a
uniqueness theorem (``given a prime factorisation $n = \prod_1^k p_i^{e_i}$, it
is unique up to ordering'') and a corollary combining existence with uniqueness.
Neither is formalised here, and their absence is not an open gap. It is recorded
in this chapter so that a reader comparing the development against the source
does not have to guess.

Two independent reasons:

\begin{itemize}
  \item \emph{It is not needed.} Nothing in the route to Euclid's theorem
    consumes uniqueness. The infinitude argument
    (Theorem~\ref{thm:not_exists_prime_list}) needs only that $m + 1$ has
    \emph{some} prime divisor, which is
    Corollary~\ref{cor:exists_prime_dvd}. The source's author reaches the same
    conclusion in the notes themselves, heading the uniqueness theorem
    % SOURCE: [euclid-notes], the heading above the second `theroem:` (read from references/euclid-infinitude-of-primes.md)
    % SOURCE QUOTE: "## Note: I dont think we actually need this uniqueness"
    \textit{``Note: I dont think we actually need this uniqueness''}.
  \item \emph{It is excluded by directive.} Uniqueness of prime factorisation is
    outside this project's stated scope, so it is not to be added and its absence
    is not to be treated as unfinished work.
\end{itemize}

The source's closing corollary --- for every $n > 1$ there exist \emph{unique}
distinct primes with positive exponents whose product is $n$ --- is therefore
established here exactly in its existence half,
Theorem~\ref{thm:exists_factorization_pow}; the uniqueness half is the excluded
result above, not an oversight. Every other named result of the source, including
each definition, is proved in this development from first principles.
